
\sectionTitle{Select Research Projects}{}
\begin{projects}

\project
  { Modeling and Control of Epidemiological Processes for HPV Spread and Cervical Cancer}{\textcolor{magenta}{NDMC IITB}}
  { \textit{Advisor: \href{https://www.iitmz.ac.in/schools/engineering-and-science/faculty/dr-nirav-bhatt}{Prof. Nirav Bhatt}} }
  {\begin{itemize}
  \setlength\itemsep{0.3em}
  \item Developed and calibrated age-structured compartmental models for HPV transmission and cervical carcinogenesis, aligned with Indian epidemiological baselines.
  \item Extended the framework to  integrate age-targeted HPV vaccination and population screening, enabling impact and evaluation of combined intervention policy scenarios.
  \item Conducted simulations quantifying the effects of varying vaccination/screening coverage, target age bands, and introduction timing on infections, precancerous lesions, cancer cases, and mortality.
  \item Generated State-Level Population Projections data using Bayesian hierarchical modeling \& Logistic-window functions 
  \item Built a decision analytic framework by solving an optimal-control problem to allocate time-varying vaccination and screening intensities under resource constraints, minimizing long-run cervical-cancer burden.
\end{itemize} }

\project
  {Predicting Global HPV and CIN Burden from Limited and Missing Data using Machine Learning}{\textcolor{magenta}{NDMC IITB}}
  { \textit{Advisor: \href{https://www.iitmz.ac.in/schools/engineering-and-science/faculty/dr-nirav-bhatt}{Prof. Nirav Bhatt}} }
  {\begin{itemize}
  \setlength\itemsep{0.3em}
  \item Curated multi-country epidemiologic,  sociodemographic \& health system covariates; engineered STI-burden and disease-incidence indices; engineered composite risk indices and applied imputation to address missingness.
  \item Built and cross-validated machine-learning models on imbalanced targets using data rebalancing using \texttt{SMOGN}.
  \item Predicted the Global and Indian state–level disease burden estimates with diagnostic analyses.
\end{itemize}}




\project
  {Agent\textendash Based Modeling with \texttt{HPVsim} for State\textendash Level HPV/Cervical\textendash Cancer Dynamics}{\textcolor{magenta}{NDMC IITB}}
  { \textit{Advisor: \href{https://www.iitmz.ac.in/schools/engineering-and-science/faculty/dr-nirav-bhatt}{Prof. Nirav Bhatt}} }
  {\begin{itemize}
  \setlength\itemsep{0.3em}
  \item Extended \href{https://www.idmod.org/tool/hpvsim/}{\texttt{HPVsim}} with maintained patches to enable subnational stratification with state\textendash level demographic and epidemiological inputs
   \item Performed state-level microsimulations for India incorporating variation in sexual behavior and demographic patterns.
\end{itemize}}
\end{projects}
\vspace{-3mm}


\begin{projects}

\project
  {Longitudinal Growth and Malnutrition Dynamics in Early Childhood: AMANHI-ACT Cohorts}
  {\textcolor{magenta}{CPHK}}
  {\textit{Advisors: \href{https://publichealth.jhu.edu/faculty/613/sunil-sazawal}{Dr.Sunil Sazawal}, \href{https://www.iitmz.ac.in/schools/engineering-and-science/faculty/dr-nirav-bhatt}{Prof. Nirav Bhatt}}}
  {\begin{itemize}\setlength\itemsep{0.5em}
      \item \textbf{SITAR Growth-Curve Modeling (Bangladesh, Pakistan, Tanzania).}
        \begin{itemize}\setlength\itemsep{0.3em}
    \item Fitted nonlinear mixed-effects \href{https://cran.r-project.org/web/packages/sitar/index.html}{\texttt{SITAR}} models for height and weight, extracting size, timing, and intensity parameters with subgroup stratification (sex, SGA/AGA, birth-weight classes).
    \item Generated mean and velocity curves and random-effects correlation maps to characterize cohort-specific growth trajectories and identify early growth faltering signals.
  \end{itemize}

      \item   \textbf{Multi-State Markov Modeling of Malnutrition Transitions.}
       \begin{itemize}\setlength\itemsep{0.3em}
    \item Defined seven nutritional states (Normal; Underweight, Stunted, Wasted; and compounded states) and estimated continuous-time transition intensities and one-year transition probabilities from cohort follow-ups.
    \item Quantified Median First Passage/Recurrence Times and lifetime time-in-state to reveal persistence of stunting and rapid cycling among compounded deficits, versus relatively faster recovery from wasting.
  \end{itemize}

  \end{itemize}}

\project
  {Cost\textendash Effectiveness Analysis using Markov Modeling for Opportunistic Prostate Cancer Screening}{\textcolor{magenta}{Apollo}}
  { \textit{Advisors: \href{https://www.apollohospitals.com/cancer-treatment-centres/cancer-hospital/chennai-apcc/radiation/dr-srinivas-chilukuri/}{Dr. Srinivas Chilukuri},  \href{https://www.iitmz.ac.in/schools/engineering-and-science/faculty/dr-nirav-bhatt}{Prof. Nirav Bhatt}} }
  {\begin{itemize}
  \setlength\itemsep{0.3em}
  \item Developed a discrete\textendash time cohort Markov model of prostate\textendash cancer natural history and screening pathways, parameterized with urban Indian epidemiology over a 20\textendash year horizon.
  \item Performed cost\textendash effectiveness analysis from a health\textendash system perspective, estimating costs and QALYs and reporting the incremental cost\textendash effectiveness ratio (ICER) versus no screening.
  \item Conducted uncertainty analysis (deterministic and probabilistic) and scenario testing to assess robustness of clinical and economic outcomes.
\end{itemize}}


  


    
\end{projects}

